\chapter{Wartość oczekiwana, momenty i podstawowe nierówności}
\label{ch:expectation}

\providecommand{\atlasref}[2]{#2}

\section{Całka względem rozkładu}
Prawdopodobieństwo odpowiada na pytanie, jak często wydarza się
zdarzenie. Wartość oczekiwana odpowiada na inne pytanie:
jaka jest przeciętna wielkość liczbowego wyniku, jeśli ważymy
możliwości ich prawdopodobieństwami? W modelu skończonym jest
to ważona suma. Całka względem miary probabilistycznej jest
jej konsekwentnym uogólnieniem na zmienne losowe, dla których
całka jest określona.
Nazwa „oczekiwana” nie oznacza wartości najbardziej prawdopodobnej;
może to być wartość, której zmienna nigdy nie przyjmuje.

\begin{definicja}
Niech $X$ będzie rzeczywistą zmienną losową na przestrzeni
probabilistycznej $(\Omega,\mathcal F,\Prob)$. Piszemy
$X^+=\max(X,0)$ oraz $X^-=\max(-X,0)$, tak że
$X=X^+-X^-$ i $|X|=X^++X^-$.
Jeśli $X\geq0$, definiujemy
$\E X:=\int_\Omega X(\omega)\,\dd\Prob(\omega)\in[0,\infty]$.
Jeżeli $X\geq0$ tylko prawie na pewno, to $\E X=\E X^+$,
gdyż $X^-$ jest wtedy zerowa prawie na pewno, a całka
nie zmienia się przy modyfikacji prawie wszędzie
(\atlasref{ME-07::thm:ae-integral-invariance}{niezmienniczość całki}).
Jeśli $\E|X|<\infty$, nazywamy $X$ całkowalną i definiujemy
$\E X:=\int_\Omega X(\omega)\,\dd\Prob(\omega)\in\R$.
W przypadku ogólnym, jeśli przynajmniej jedna z liczb $\E X^+,\E X^-$ jest
skończona, dopuszczamy rozszerzoną wartość oczekiwaną
$\E X=\E X^+-\E X^-$. Gdy obie są nieskończone,
wartość oczekiwana nie jest określona.
Jest to definicja całki Lebesgue'a względem $\Prob$;
por. \atlasref{ME-07::thm:integrability-criterion}{kryterium całkowalności}.
Dla zmiennej zespolonej $Z$ o $\E|Z|<\infty$ definiujemy
$\E Z=\E\operatorname{Re}Z+i\E\operatorname{Im}Z$.
Obie części są całkowalne, ponieważ ich moduły nie przekraczają
$|Z|$; por. \atlasref{ME-08::def:complex-integral}{definicję całki zespolonej}.
\end{definicja}

Dla całkowalnych rzeczywistych zmiennych losowych $X,Y$
na tej samej przestrzeni probabilistycznej oraz $a,b\in\R$ zmienna $aX+bY$
jest całkowalna, gdyż $|aX+bY|\leq |a||X|+|b||Y|$.
Z \atlasref{ME-08::thm:integral-linear-real}{liniowości całki}
wynika $\E(aX+bY)=a\E X+b\E Y$.
Ponadto dla stałej $c\in\R$ mamy $\E c=c\Prob(\Omega)=c$,
a z monotoniczności całki, jeśli $X\leq Y$ prawie na pewno,
to $\E X\leq\E Y$. Stosujemy tu monotoniczność do
nieujemnych części lub, równoważnie, do różnicy $Y-X$,
której całka jest nieujemna. W szczególności
$|\E X|\leq\E|X|$, ponieważ $-|X|\leq X\leq |X|$.
Dla całkowalnych zmiennych zespolonych obowiązuje liniowość
ze współczynnikami zespolonymi oraz $|\E Z|\leq\E|Z|$
(\atlasref{ME-08::thm:integral-complex-linear}{liniowość i oszacowanie modułu całki zespolonej}).
Z definicji całki zespolonej wynika również
$\E\overline Z=\overline{\E Z}$.

Dla nieujemnej zmiennej losowej $U$ zachodzi $\E U=0$
wtedy i tylko wtedy, gdy $U=0$ prawie na pewno.
Jeśli $\E U=0$, to dla każdego $n\geq1$ mamy
$0\leq n^{-1}\chi_{\{U\geq 1/n\}}\leq U$.
Monotoniczność całki nieujemnej daje
$n^{-1}\Prob(U\geq1/n)\leq\E U=0$.
Zatem $\{U>0\}=\bigcup_{n\geq1}\{U\geq1/n\}$ ma
prawdopodobieństwo zero. Odwrotna implikacja wynika
z niezmienniczości całki względem równości prawie wszędzie.
Własność tę stosujemy również do zmiennych nieujemnych
prawie na pewno, zastępując je ich częściami dodatnimi.

\begin{twierdzenie}[Przenoszenie całkowania na rozkład]
Niech $X$ będzie rzeczywistą zmienną losową o rozkładzie $\mu_X$,
a $h:\R\to\mathbb C$ funkcją borelowską (w szczególności
może być rzeczywista). Jeżeli $h$ przyjmuje wyłącznie
wartości nieujemne rzeczywiste, to $h(X)$ jest nieujemna.
Dla dowolnej funkcji borelowskiej $h:\R\to\mathbb C$
zachodzi
\[
\E|h(X)|=\int_\R|h(x)|\,\dd\mu_X(x)\in[0,\infty].
\]
W szczególności $h(X)$ jest całkowalna wtedy i tylko wtedy,
gdy $h$ jest całkowalna względem $\mu_X$.
Jeżeli $h$ jest rzeczywista i nieujemna lub $h$ jest
całkowalna względem $\mu_X$, to
\begin{equation}\label{eq:lotus}
\E[h(X)]=\int_\R h(x)\,\dd\mu_X(x).
\end{equation}
W szczególności, dla dyskretnego $X$ o masach $p(x)$ prawa strona
jest sumą $\sum_{x\in D}h(x)p(x)$, gdzie
$D=\{x:p(x)>0\}$ jest co najwyżej przeliczalny.
Dla $X$ o gęstości $f_X$ (wybieramy wersję wszędzie
skończoną) prawa strona jest całką
$\int_\R h(x)f_X(x)\,\dd\lambda_1(x)$.
W przypadku funkcji rzeczywistych ze znakiem lub zespolonych
obowiązuje tu założenie całkowalności; w przypadku nieujemnym
dopuszczamy wartość nieskończoną.
\end{twierdzenie}
\begin{proof}
Jest to szczególny przypadek
\atlasref{ME-13::thm:pushforward-integral-13}{twierdzenia o całkowaniu względem miary obrazowej}.
Podamy dowód w notacji probabilistycznej. Złożenie $h\circ X$
jest mierzalne, ponieważ $X$ i $h$ są mierzalne.

Dla funkcji charakterystycznej $h=\chi_B$, gdzie
$B\in\mathcal B(\R)$, mamy $h(X)=\chi_{X^{-1}(B)}$ oraz
$\E[h(X)]=\Prob(X\in B)=\mu_X(B)
=\int_\R\chi_B(x)\,\dd\mu_X(x)$.
Dla nieujemnej funkcji prostej
$h=\sum_{j=1}^r a_j\chi_{B_j}$, gdzie $a_j\geq0$,
a $B_j$ są zbiorami borelowskimi,
równość wynika z liniowości całki i poprzedniego przypadku.

Dla dowolnej nieujemnej funkcji mierzalnej $h$
wybieramy ciąg nieujemnych borelowskich funkcji prostych
$h_n\uparrow h$
(\atlasref{ME-06::thm:simple-up}{aproksymacja funkcjami prostymi}).
Wówczas $h_n(X)\uparrow h(X)$ na $\Omega$.
Stosując \atlasref{ME-08::thm:mct}{twierdzenie o zbieżności monotonicznej}
po obu stronach już udowodnionych równości, dostajemy
\[
\E[h(X)]=\lim_n\E[h_n(X)]
=\lim_n\int_\R h_n(x)\,\dd\mu_X(x)
=\int_\R h(x)\,\dd\mu_X(x).
\]
Jeśli $h$ jest rzeczywista i całkowalna względem $\mu_X$,
rozkładamy $h=h^+-h^-$ i stosujemy pierwszy przypadek
do obu nieujemnych części. Stosując również przypadek
nieujemny do $|h|$, otrzymujemy
$\E|h(X)|=\int_\R|h(x)|\,\dd\mu_X(x)<\infty$.
Obie odpowiadające całki są zatem skończone i można
odjąć uzyskane równości.

Dla zespolonej funkcji całkowalnej $h$ stosujemy wynik
do $\operatorname{Re}h$ i $\operatorname{Im}h$, które są
całkowalne, a następnie korzystamy z definicji całki zespolonej.
Równość dla $|h|$ została już udowodniona w przypadku
nieujemnym, również wtedy, gdy obie strony są nieskończone.

Dla rozkładu dyskretnego $\mu_X(D)=1$. Dla $h\geq0$
rozkładamy całkę na sumę po singletonach $\{x\}$,
$x\in D$, stosując
\atlasref{ME-08::cor:integral-countable-partition}{addytywność całki względem zbioru całkowania}.
Całka po singletonie wynosi $h(x)p(x)$, a po
$\R\setminus D$ jest zerowa. Dla całkowalnego rzeczywistego $h$
stosujemy to do $h^+$ i $h^-$; szereg jest bezwzględnie
zbieżny, bo $\sum_{x\in D}|h(x)|p(x)<\infty$.
Dla funkcji zespolonej stosujemy wzór do obu części,
a bezwzględna zbieżność wynika z tego samego oszacowania.
W przypadku gęstości stosujemy
\atlasref{ME-18::lem:integration-density-18}{wzór na całkę względem miary określonej gęstością}
do $\mu_X=f_X\lambda_1$ (obie miary na $\mathcal B(\R)$).
Możemy wybrać wszędzie skończoną wersję $f_X$, ponieważ
$\int_\R f_X(x)\,\dd\lambda_1(x)=1$, a zbiór
$\{f_X=\infty\}$ jest borelowski i ma miarę zero.
Dla funkcji nieujemnej wzór obowiązuje także przy
nieskończonej całce; dla całkowalnej funkcji rzeczywistej
lub zespolonej całkowalność iloczynu $hf_X$ wynika z
$\int_\R |h(x)|f_X(x)\,\dd\lambda_1(x)=\int_\R |h(x)|\,\dd\mu_X(x)$.
\end{proof}


\begin{przyklad}[Średnia nie musi być możliwym wynikiem]
Niech $X$ przyjmuje wartości $0$ i $1$ z prawdopodobieństwami
po $1/2$. Wówczas $\E X=1/2$, ale $\Prob(X=1/2)=0$.
Średnia jest środkiem ciężkości rozkładu, a nie typową
zrealizowaną wartością. Dla bardzo nierównomiernych rozkładów
może ona znajdować się daleko od większości realizacji.
\end{przyklad}

\begin{definicja}
Dla całkowalnej rzeczywistej zmiennej losowej $X$
definiujemy średnią $m_X:=\E X$. Jeżeli $\E X^2<\infty$, to $X$ i funkcja stała $1$
należą do $L^2(\Prob)$. Z nierówności
Cauchy'ego--Schwarza dla tych funkcji mamy
$\E|X|\leq(\E X^2)^{1/2}$, gdyż $\Prob(\Omega)=1$.
Definiujemy wariancję $\Var(X):=\E[(X-\E X)^2]$.
Jest ona skończona, ponieważ
$(X-\E X)^2\leq 2X^2+2(\E X)^2$.
Odchyleniem standardowym nazywamy $\sqrt{\Var(X)}$.
Wariancja jest nieujemna. Jest równa zero wtedy i tylko wtedy,
gdy $X=\E X$ prawie na pewno: całka nieujemnej funkcji
$(X-\E X)^2$ jest zerowa wtedy i tylko wtedy, gdy ta funkcja
znika prawie wszędzie.
Dla rzeczywistych zmiennych $X,Y$ na tej samej przestrzeni
probabilistycznej o skończonych drugich momentach definiujemy
kowariancję $\Cov(X,Y):=\E[(X-\E X)(Y-\E Y)]$.
Z nierówności Cauchy'ego--Schwarza wynika
$\E|(X-\E X)(Y-\E Y)|\leq
\sqrt{\Var(X)\Var(Y)}<\infty$, więc kowariancja jest
dobrze określona. Użyta nierówność to przypadek $p=q=2$
\atlasref{ME-14::thm:holder-14}{nierówności Höldera}.
\end{definicja}

\begin{stwierdzenie}[Algebra momentów]
Dla rzeczywistych zmiennych losowych $X,Y$ na wspólnej
przestrzeni probabilistycznej, o skończonych drugich
momentach, oraz $a,b\in\R$ zachodzą
\begin{align*}
\Var(X)&=\E X^2-(\E X)^2, &
\Var(aX+b)&=a^2\Var(X),\\
\Cov(X,Y)&=\E[XY]-\E X\,\E Y, &
\Var(X+Y)&=\Var(X)+\Var(Y)+2\Cov(X,Y).
\end{align*}
Jeżeli $X$ i $Y$ są niezależne, to $\E[XY]=\E X\,\E Y$,
a zatem $\Cov(X,Y)=0$ i $\Var(X+Y)=\Var(X)+\Var(Y)$.
\end{stwierdzenie}
\begin{proof}
Wszystkie podane wyrażenia są skończone,
ponieważ nierówność Cauchy'ego--Schwarza zapewnia
$\E|XY|\leq(\E X^2)^{1/2}(\E Y^2)^{1/2}<\infty$.
Ponadto $(aX+b)^2\leq2a^2X^2+2b^2$ i
$(X+Y)^2\leq2X^2+2Y^2$, więc również przekształcenia
w tezie mają skończone drugie momenty.
Możemy zatem korzystać z
\atlasref{ME-08::thm:integral-linear-real}{liniowości całki funkcji całkowalnych}.
Oznaczając $m=\E X$, rozwijamy
$\E(X-m)^2=\E X^2-2m\E X+m^2=\E X^2-m^2$.
Ponadto $\E(aX+b)=a\E X+b$, więc
$aX+b-\E(aX+b)=a(X-\E X)$,
a stąd $\Var(aX+b)=a^2\Var(X)$.

Dla $m_X=\E X$, $m_Y=\E Y$ mamy
\[
\begin{aligned}
\Cov(X,Y)
&=\E[(X-m_X)(Y-m_Y)]\\
&=\E[XY]-m_Y\E X-m_X\E Y+m_Xm_Y\\
&=\E[XY]-m_Xm_Y.
\end{aligned}
\]
Ponieważ
$X+Y-\E(X+Y)=(X-m_X)+(Y-m_Y)$,
rozwinięcie kwadratu i całkowanie daje
$\Var(X+Y)=\Var(X)+\Var(Y)+2\Cov(X,Y)$.

Jeśli $X,Y$ są niezależne, ich rozkład łączny to
$\mu_X\otimes\mu_Y$ (rozdział~\ref{ch:random-variables}).
Stosując \atlasref{ME-13::thm:pushforward-integral-13}{wzór na całkę względem miary obrazowej}
do nieujemnej funkcji $(x,y)\mapsto|xy|$, otrzymujemy
$\int_{\R^2}|xy|\,\dd(\mu_X\otimes\mu_Y)(x,y)
=\E|XY|<\infty$. Spełniona jest zatem hipoteza całkowalności
funkcji $(x,y)\mapsto xy$ w twierdzeniu Fubiniego. Z \atlasref{ME-12::thm:fubini-12}{twierdzenia Fubiniego}
i wzoru na całkę względem miary obrazowej otrzymujemy
\[
\E[XY]=\int_{\R^2}xy\,\dd(\mu_X\otimes\mu_Y)(x,y)
=\left(\int_\R x\,\dd\mu_X(x)\right)
 \left(\int_\R y\,\dd\mu_Y(y)\right)
=\E X\,\E Y.
\]
To daje zerową kowariancję i wzór na wariancję sumy.
Odwrotna implikacja nie jest prawdziwa i nie wolno
jej wykorzystywać bez dodatkowych założeń.
\end{proof}

\begin{uwaga}
Zerowa kowariancja nie oznacza na ogół niezależności.
Jeśli $X$ jest jednostajna na $[-1,1]$, a $Y=X^2$, to
obie zmienne są ograniczone prawie na pewno, zatem mają
skończone drugie momenty. Stąd
$\Cov(X,Y)=\E X^3-\E X\,\E X^2=0$,
ponieważ $\E X=\frac12\int_{-1}^1 x\,\dd\lambda_1(x)=0$
oraz $\E X^3=\frac12\int_{-1}^1 x^3\,\dd\lambda_1(x)=0$
z nieparzystości funkcji podcałkowych.
Jednak dla $A=\{|X|\leq1/2\}$ i $B=\{Y\leq1/4\}$
mamy $A=B$, a z rozkładu jednostajnego
$\Prob(A)=\Prob(B)=1/2$ i $\Prob(A\cap B)=1/2\neq1/4$.
Zatem zmienne nie są niezależne.
\end{uwaga}

\begin{stwierdzenie}[Momenty funkcji charakterystycznej zdarzenia]
Niech $(\Omega,\mathcal F,\Prob)$ będzie przestrzenią
probabilistyczną. Jeżeli $A\in\mathcal F$ i $I=\chi_A$, to
$\E I=\Prob(A)$ oraz
$\Var(I)=\Prob(A)(1-\Prob(A))$.
Dla dowolnych zdarzeń $A,B\in\mathcal F$ zachodzi
$\Cov(\chi_A,\chi_B)=\Prob(A\cap B)-\Prob(A)\Prob(B)$.
\end{stwierdzenie}
\begin{proof}
Funkcja $I$ jest mierzalna, ograniczona i spełnia $I^2=I$.
Z definicji całki funkcji charakterystycznej
$\E I=\int_\Omega\chi_A(\omega)\,\dd\Prob(\omega)=\Prob(A)$.
Ze wzoru na wariancję otrzymujemy
$\Var(I)=\E I^2-(\E I)^2=\Prob(A)-\Prob(A)^2$.
Ponadto $\chi_A\chi_B=\chi_{A\cap B}$, więc ze wzoru
na kowariancję wynika
$\Cov(\chi_A,\chi_B)=\E\chi_{A\cap B}-\E\chi_A\,\E\chi_B
=\Prob(A\cap B)-\Prob(A)\Prob(B)$.
\end{proof}

\section{Nierówności Markowa i Czebyszewa}
Nie zawsze umiemy wyznaczyć cały rozkład, ale czasem znamy
średnią albo wariancję. Czy można na tej podstawie oszacować
prawdopodobieństwo dużego wyniku lub dużego odchylenia od średniej?
Nierówność Markowa opiera się tylko na nieujemności, a nierówność
Czebyszewa stosuje tę samą ideę do kwadratu odchylenia.
Ich siła tkwi w uniwersalności: nie wymagają normalności,
jednostajności ani nawet znajomości gęstości.

\begin{twierdzenie}[Nierówność Markowa]
Niech $X$ będzie rzeczywistą zmienną losową na przestrzeni
probabilistycznej $(\Omega,\mathcal F,\Prob)$. Jeżeli
$X\geq0$ prawie na pewno i $a>0$, to
$\Prob(X\geq a)\leq\E X/a$, gdzie dopuszczamy $\E X=\infty$.
Jeżeli $\E X<\infty$, równość zachodzi wtedy i tylko wtedy,
gdy $X\in\{0,a\}$ prawie na pewno. Stała $1$ jest optymalna.
\end{twierdzenie}
\begin{proof}
Zdarzenie $A=\{X\geq a\}$ jest mierzalne. Na $A$ zachodzi
$X\geq a$, a na $\Omega\setminus A$ mamy $a\chi_A=0$.
Zatem $a\chi_A\leq X^+$ na całej przestrzeni.
Ponieważ $X^-=0$ prawie na pewno, zachodzi $\E X=\E X^+$.
Z monotoniczności całki nieujemnej otrzymujemy
$a\Prob(A)=\E[a\chi_A]\leq\E X^+=\E X$.
Jest to probabilistyczna postać
\atlasref{ME-07::thm:markov-nonnegative}{nierówności Markowa dla całki nieujemnej}.
Dzielenie przez $a>0$ daje tezę. Nierówność pozostaje
prawdziwa, choć może być bezużyteczna, gdy $\E X=\infty$.
Jeżeli $\E X<\infty$, równość zachodzi wtedy i tylko wtedy,
gdy $X\in\{0,a\}$ prawie na pewno. Istotnie, różnica
$X^+-a\chi_A$ jest nieujemna, a jej całka jest zerowa
dokładnie wtedy, gdy różnica znika prawie wszędzie.
Na $A$ oznacza to $X=a$, a poza $A$ oznacza $X^+=0$,
czyli $X=0$ poza zbiorem zerowym. Odwrotna implikacja
wynika przez podstawienie. Biorąc $X=a$ z prawdopodobieństwem
jeden, otrzymujemy równość z dodatnią prawą stroną;
stałej $1$ w nierówności nie można więc zmniejszyć.
\end{proof}

\begin{twierdzenie}[Nierówność Czebyszewa]
Niech $X$ będzie rzeczywistą zmienną losową o skończonym
drugim momencie. Wtedy dla każdego $\varepsilon>0$ zachodzi
\[
\Prob(|X-\E X|\geq\varepsilon)
\leq\frac{\Var(X)}{\varepsilon^2}.
\]
Równość zachodzi wtedy i tylko wtedy, gdy
$|X-\E X|\in\{0,\varepsilon\}$ prawie na pewno.
Stała $1$ jest optymalna.
\end{twierdzenie}
\begin{proof}
Z założenia o drugim momencie wynika całkowalność $X$,
więc $m=\E X$ jest skończoną liczbą rzeczywistą.
Wprowadźmy nieujemną zmienną
$Y=(X-m)^2$, mierzalną jako złożenie $X$ z funkcją ciągłą. Jej wartość oczekiwana jest równa
$\E Y=\Var(X)<\infty$. Ponadto dla $\varepsilon>0$
zdarzenia
$\{Y\geq\varepsilon^2\}$ oraz
$\{|X-\E X|\geq\varepsilon\}$ są identyczne.
Stosując właśnie udowodnioną nierówność Markowa
do zmiennej $Y$ i progu $\varepsilon^2$, otrzymujemy
\[
\Prob(|X-\E X|\geq\varepsilon)
=\Prob(Y\geq\varepsilon^2)
\leq\frac{\E Y}{\varepsilon^2}
=\frac{\Var(X)}{\varepsilon^2}.
\]
Nie używaliśmy żadnych informacji o kształcie rozkładu $X$.
Z warunku równości w nierówności Markowa wynika, że
równość zachodzi dokładnie wtedy, gdy
$|X-\E X|\in\{0,\varepsilon\}$ prawie na pewno.
W szczególności zmienna przyjmująca wartości
$-\varepsilon$ i $\varepsilon$ z prawdopodobieństwami
po $1/2$ ma średnią zero, wariancję $\varepsilon^2$
i osiąga równość. Stała $1$ jest optymalna.
\end{proof}

\begin{przyklad}[Kontrola odchylenia sumy]
Niech $X_1,\ldots,X_n$ będą niezależnymi rzeczywistymi
zmiennymi losowymi na wspólnej przestrzeni probabilistycznej,
o wspólnej średniej $m$ i wariancji $\sigma^2<\infty$,
gdzie $n\geq1$ jest liczbą całkowitą.
Dla średniej $\overline X_n:=n^{-1}\sum_{i=1}^nX_i$
liniowość wartości oczekiwanej daje
$\E\overline X_n=n^{-1}\sum_{i=1}^n m=m$.
Skończoność drugiego momentu sumy wynika z nierówności
$(\sum_{i=1}^n X_i)^2\leq n\sum_{i=1}^n X_i^2$,
która jest nierównością Cauchy'ego--Schwarza dla wektorów
$(X_1(\omega),\ldots,X_n(\omega))$ oraz $(1,\ldots,1)$.
Rozwijając kwadrat wycentrowanej sumy, otrzymujemy
\[
\Var(\overline X_n)=\frac{1}{n^2}
\left(\sum_{i=1}^n\Var(X_i)+
2\sum_{1\leq i<j\leq n}\Cov(X_i,X_j)\right).
\]
Kowariancje dla $i\ne j$ są zerowe z niezależności
i poprzedniego stwierdzenia o algebrze momentów.
Stąd $\Var(\overline X_n)=n^{-2}n\sigma^2=\sigma^2/n$.
Z nierówności Czebyszewa dla każdego $\varepsilon>0$ wynika
$\Prob(|\overline X_n-m|\geq\varepsilon)
\leq\sigma^2/(n\varepsilon^2)$.
Jest to zalążek słabego prawa wielkich liczb, którego
pełne sformułowania i uogólnienia pojawią się w rozdziale~\ref{ch:lln}.
Rysunek~\ref{fig:rp1-05-chebyshev-average} przedstawia zależność oszacowania od liczby obserwacji i progu odchylenia.
\end{przyklad}

\begin{figure}[htbp]
\centering
\begin{tikzpicture}[x=1.8cm,y=4.3cm,>=Stealth]
\draw[->] (0,0)--(5.3,0) node[right] {$\varepsilon/\sigma$};
\draw[->] (0,0)--(0,1.13) node[above] {ograniczenie};
\foreach \x in {1,2,3,4,5} \node[below,font=\small] at (\x,0) {$\x$};
\foreach \y/\lab in {.25/{1/4},.5/{1/2},.75/{3/4},1/1} {
 \draw[gray!20] (0,\y)--(5,\y);
 \node[left,font=\small] at (0,\y) {$\lab$};
}
\draw[blue!75!black,very thick] plot coordinates {(0,1) (0.05,1) (0.1,1) (0.15,1) (0.2,1) (0.25,1) (0.3,1) (0.35,1) (0.4,1) (0.45,1) (0.5,1) (0.55,1) (0.6,1) (0.65,1) (0.7,1) (0.75,1) (0.8,1) (0.85,1) (0.9,1) (0.95,1) (1,1) (1.05,0.9070295) (1.1,0.8264463) (1.15,0.7561437) (1.2,0.6944444) (1.25,0.64) (1.3,0.591716) (1.35,0.5486968) (1.4,0.5102041) (1.45,0.4756243) (1.5,0.4444444) (1.55,0.4162331) (1.6,0.390625) (1.65,0.3673095) (1.7,0.3460208) (1.75,0.3265306) (1.8,0.308642) (1.85,0.2921841) (1.9,0.2770083) (1.95,0.2629849) (2,0.25) (2.05,0.2379536) (2.1,0.2267574) (2.15,0.2163332) (2.2,0.2066116) (2.25,0.1975309) (2.3,0.1890359) (2.35,0.1810774) (2.4,0.1736111) (2.45,0.1665973) (2.5,0.16) (2.55,0.153787) (2.6,0.147929) (2.65,0.1423994) (2.7,0.1371742) (2.75,0.1322314) (2.8,0.127551) (2.85,0.1231148) (2.9,0.1189061) (2.95,0.1149095) (3,0.1111111) (3.05,0.107498) (3.1,0.1040583) (3.15,0.1007811) (3.2,0.0976562) (3.25,0.0946746) (3.3,0.0918274) (3.35,0.0891067) (3.4,0.0865052) (3.45,0.084016) (3.5,0.0816327) (3.55,0.0793493) (3.6,0.0771605) (3.65,0.075061) (3.7,0.073046) (3.75,0.0711111) (3.8,0.0692521) (3.85,0.067465) (3.9,0.0657462) (3.95,0.0640923) (4,0.0625) (4.05,0.0609663) (4.1,0.0594884) (4.15,0.0580636) (4.2,0.0566893) (4.25,0.0553633) (4.3,0.0540833) (4.35,0.0528471) (4.4,0.0516529) (4.45,0.0504987) (4.5,0.0493827) (4.55,0.0483033) (4.6,0.047259) (4.65,0.0462481) (4.7,0.0452694) (4.75,0.0443213) (4.8,0.0434028) (4.85,0.0425125) (4.9,0.0416493) (4.95,0.0408122) (5,0.04)};
\draw[orange!85!black,dashed,very thick] plot coordinates {(0,1) (0.05,1) (0.1,1) (0.15,1) (0.2,1) (0.25,1) (0.3,1) (0.35,1) (0.4,1) (0.45,1) (0.5,1) (0.55,0.8264463) (0.6,0.6944444) (0.65,0.591716) (0.7,0.5102041) (0.75,0.4444444) (0.8,0.390625) (0.85,0.3460208) (0.9,0.308642) (0.95,0.2770083) (1,0.25) (1.05,0.2267574) (1.1,0.2066116) (1.15,0.1890359) (1.2,0.1736111) (1.25,0.16) (1.3,0.147929) (1.35,0.1371742) (1.4,0.127551) (1.45,0.1189061) (1.5,0.1111111) (1.55,0.1040583) (1.6,0.0976562) (1.65,0.0918274) (1.7,0.0865052) (1.75,0.0816327) (1.8,0.0771605) (1.85,0.073046) (1.9,0.0692521) (1.95,0.0657462) (2,0.0625) (2.05,0.0594884) (2.1,0.0566893) (2.15,0.0540833) (2.2,0.0516529) (2.25,0.0493827) (2.3,0.047259) (2.35,0.0452694) (2.4,0.0434028) (2.45,0.0416493) (2.5,0.04) (2.55,0.0384468) (2.6,0.0369822) (2.65,0.0355999) (2.7,0.0342936) (2.75,0.0330579) (2.8,0.0318878) (2.85,0.0307787) (2.9,0.0297265) (2.95,0.0287274) (3,0.0277778) (3.05,0.0268745) (3.1,0.0260146) (3.15,0.0251953) (3.2,0.0244141) (3.25,0.0236686) (3.3,0.0229568) (3.35,0.0222767) (3.4,0.0216263) (3.45,0.021004) (3.5,0.0204082) (3.55,0.0198373) (3.6,0.0192901) (3.65,0.0187652) (3.7,0.0182615) (3.75,0.0177778) (3.8,0.017313) (3.85,0.0168663) (3.9,0.0164366) (3.95,0.0160231) (4,0.015625) (4.05,0.0152416) (4.1,0.0148721) (4.15,0.0145159) (4.2,0.0141723) (4.25,0.0138408) (4.3,0.0135208) (4.35,0.0132118) (4.4,0.0129132) (4.45,0.0126247) (4.5,0.0123457) (4.55,0.0120758) (4.6,0.0118147) (4.65,0.011562) (4.7,0.0113173) (4.75,0.0110803) (4.8,0.0108507) (4.85,0.0106281) (4.9,0.0104123) (4.95,0.010203) (5,0.01)};
\draw[teal!75!black,dashdotted,very thick] plot coordinates {(0,1) (0.05,1) (0.1,1) (0.15,1) (0.2,1) (0.25,1) (0.3,0.6944444) (0.35,0.5102041) (0.4,0.390625) (0.45,0.308642) (0.5,0.25) (0.55,0.2066116) (0.6,0.1736111) (0.65,0.147929) (0.7,0.127551) (0.75,0.1111111) (0.8,0.0976562) (0.85,0.0865052) (0.9,0.0771605) (0.95,0.0692521) (1,0.0625) (1.05,0.0566893) (1.1,0.0516529) (1.15,0.047259) (1.2,0.0434028) (1.25,0.04) (1.3,0.0369822) (1.35,0.0342936) (1.4,0.0318878) (1.45,0.0297265) (1.5,0.0277778) (1.55,0.0260146) (1.6,0.0244141) (1.65,0.0229568) (1.7,0.0216263) (1.75,0.0204082) (1.8,0.0192901) (1.85,0.0182615) (1.9,0.017313) (1.95,0.0164366) (2,0.015625) (2.05,0.0148721) (2.1,0.0141723) (2.15,0.0135208) (2.2,0.0129132) (2.25,0.0123457) (2.3,0.0118147) (2.35,0.0113173) (2.4,0.0108507) (2.45,0.0104123) (2.5,0.01) (2.55,0.0096117) (2.6,0.0092456) (2.65,0.0089) (2.7,0.0085734) (2.75,0.0082645) (2.8,0.0079719) (2.85,0.0076947) (2.9,0.0074316) (2.95,0.0071818) (3,0.0069444) (3.05,0.0067186) (3.1,0.0065036) (3.15,0.0062988) (3.2,0.0061035) (3.25,0.0059172) (3.3,0.0057392) (3.35,0.0055692) (3.4,0.0054066) (3.45,0.005251) (3.5,0.005102) (3.55,0.0049593) (3.6,0.0048225) (3.65,0.0046913) (3.7,0.0045654) (3.75,0.0044444) (3.8,0.0043283) (3.85,0.0042166) (3.9,0.0041091) (3.95,0.0040058) (4,0.0039063) (4.05,0.0038104) (4.1,0.003718) (4.15,0.003629) (4.2,0.0035431) (4.25,0.0034602) (4.3,0.0033802) (4.35,0.0033029) (4.4,0.0032283) (4.45,0.0031562) (4.5,0.0030864) (4.55,0.003019) (4.6,0.0029537) (4.65,0.0028905) (4.7,0.0028293) (4.75,0.0027701) (4.8,0.0027127) (4.85,0.002657) (4.9,0.0026031) (4.95,0.0025508) (5,0.0025)};
\draw[blue!75!black,very thick] (3.1,.89)--(3.55,.89) node[right,font=\small] {$n=1$};
\draw[orange!85!black,very thick,dashed] (3.1,.78)--(3.55,.78) node[right,font=\small] {$n=4$};
\draw[teal!75!black,very thick,dashdotted] (3.1,.67)--(3.55,.67) node[right,font=\small] {$n=16$};
\end{tikzpicture}
\caption{Ograniczenia Czebyszewa $\min(1,1/(na^2))$,
gdzie $a=\varepsilon/\sigma>0$ i $\sigma>0$.
Wartość w zerze oznacza ciągłe przedłużenie ograniczenia.
Krzywe nie są dokładnymi prawdopodobieństwami;
ich obniżanie przy wzroście $n$ ilustruje kontrolę średniej.}
\label{fig:rp1-05-chebyshev-average}
\end{figure}


\section{Momenty wyższych rzędów i sens funkcji tworzących}

Dwa pierwsze momenty często nie wystarczają, aby odróżnić
rozkłady. Zmienna może mieć tę samą średnią i wariancję co inna,
a jednocześnie zupełnie inny kształt ogonów. Chcielibyśmy zatem
przechowywać informacje o całej rodzinie momentów, a najlepiej
w jednej funkcji, na której można wykonywać zwykłe operacje
analityczne. Dwa użyteczne pomysły polegają na obliczaniu
$\E z^X$ dla zmiennych o wartościach w $\N_0=\{0,1,2,\ldots\}$ i $\E e^{tX}$
dla zmiennych rzeczywistych. Trzeci --- wykorzystujący
$e^{itX}$ --- pozwala uniknąć problemu rozbieżności.

\begin{definicja}
Niech $X$ będzie rzeczywistą zmienną losową, a $r\geq1$
liczbą całkowitą. Jeśli $\E|X|^r<\infty$, to
$m_r(X):=\E X^r$ nazywamy momentem zwykłym rzędu $r$.
Ponieważ $|x|\leq1+|x|^r$, założenie to zapewnia
całkowalność $X$. Ponadto nierówność trójkąta i wypukłość funkcji
$t\mapsto t^r$ na $[0,\infty)$ zastosowana z wagami
$1/2,1/2$ dają
\[
|X-\E X|^r\leq 2^{r-1}(|X|^r+|\E X|^r),
\]
więc $\mu_r(X):=\E[(X-\E X)^r]$ jest skończona.
Nazywamy ją momentem centralnym rzędu $r$.
Dla dowolnego rzeczywistego $s>0$ wielkość
$\E|X|^s\in[0,\infty]$ nazywamy momentem bezwzględnym
rzędu $s$.
\end{definicja}

\begin{uwaga}
Dla nieparzystego $r$ całka $\E X^r$ może nie istnieć,
nawet gdy symetria rozkładu formalnie sugeruje zero.
Istnienie momentu zwykłego jako skończonej liczby
w sensie całki Lebesgue'a jest równoważne warunkowi
$\E|X|^r<\infty$, ponieważ $|X^r|=|X|^r$.
Na przykład dla standardowego rozkładu Cauchy'ego
(por. \eqref{eq:cauchy} w rozdziale~\ref{ch:distributions})
obie części całki definiującej średnią są nieskończone:
dla części dodatniej
$\int_0^R x/(\pi(1+x^2))\,\dd\lambda_1(x)
=\log(1+R^2)/(2\pi)\to\infty$, a część ujemna
ma tę samą całkę z symetrii gęstości.
Gdy $r>s>0$, dla $|x|\leq1$ mamy $|x|^s\leq1$,
a dla $|x|>1$ mamy $|x|^s\leq|x|^r$.
W konsekwencji $\E|X|^s\leq1+\E|X|^r$,
więc skończoność momentu bezwzględnego rzędu $r$
zapewnia skończoność wszystkich momentów bezwzględnych
dodatnich rzędów mniejszych od $r$.
\end{uwaga}

\section{Funkcja tworząca prawdopodobieństwa}

Dla zmiennej $N$ przyjmującej wartości w $\N_0:=\{0,1,2,\ldots\}$
rozkład jest ciągiem $p_k=\Prob(N=k)$. Możemy potraktować
$p_k$ jako współczynniki szeregu potęgowego. Wówczas
mnożenie funkcji tworzących odpowiada dodawaniu niezależnych
liczb wystąpień --- dokładnie tak, jak splot ciągów mas
opisuje rozkład ich sumy.

\begin{definicja}
Niech $N:\Omega\to\N_0$ będzie zmienną losową na przestrzeni
probabilistycznej $(\Omega,\mathcal F,\Prob)$.
Jej funkcją tworzącą prawdopodobieństwa nazywamy
$G_N(z):=\E z^N=\sum_{k=0}^\infty \Prob(N=k)z^k$
dla $z\in\mathbb C$, $|z|\leq1$, z konwencją $0^0=1$.
Funkcja $\omega\mapsto z^{N(\omega)}$ jest mierzalna:
przeciwobraz dowolnego zbioru borelowskiego w $\mathbb C$
jest przeliczalną sumą odpowiednich zdarzeń $\{N=k\}$. Jej moduł
nie przekracza jedności, więc funkcja $z^N$ jest całkowalna.
Równość z szeregiem wynika z \eqref{eq:lotus},
zastosowanego do funkcji borelowskiej $h$ określonej
przez $h(k)=z^k$ dla $k\in\N_0$ i $h(x)=0$ poza $\N_0$.
W ten sposób nie trzeba definiować potęgi zespolonej
$z^x$ dla dowolnego rzeczywistego $x$.
Szereg jest bezwzględnie i jednostajnie zbieżny na domkniętym
kole jednostkowym, ponieważ $|p_kz^k|\leq p_k$ i
$\sum_{k=0}^\infty p_k=1$. W szczególności $G_N$ jest
ciągła na tym kole.
\end{definicja}

\begin{twierdzenie}[Odzyskiwanie rozkładu i momentów silniowych]
Niech $N$ będzie zmienną losową o wartościach w $\N_0$,
a $G_N$ jej funkcją tworzącą. Funkcja $G_N$ jest holomorficzna
w otwartym kole jednostkowym, $G_N(1)=1$, a dla $k\in\N_0$
zachodzi $\Prob(N=k)=G_N^{(k)}(0)/k!$, gdzie
$G_N^{(0)}=G_N$. W konsekwencji funkcja $G_N$ wyznacza
rozkład $N$ jednoznacznie.
Dla całkowitego $r\geq1$ i $0\leq z<1$ mamy
\[
G_N^{(r)}(z)=\E\bigl[(N)_r z^{N-r}\chi_{\{N\geq r\}}\bigr],
\qquad (n)_r:=n(n-1)\cdots(n-r+1).
\]
W powyższym wzorze całe wyrażenie pod wartością oczekiwaną
ma wartość $0$ na $\{N<r\}$, bez obliczania $z^{N-r}$;
dla $z=0$ i $N=r$ przyjmujemy $0^0=1$.
Dla $n<r$ iloczyn $(n)_r$ jest zerowy.
W szczególności istnieje granica lewostronna
$G_N^{(r)}(1-):=\lim_{z\uparrow1}G_N^{(r)}(z)
=\E(N)_r\in[0,\infty]$.
Nie utożsamiamy tej granicy z pochodną zespoloną w punkcie
$1$, który może leżeć na brzegu koła zbieżności.
Dla $r=1$ w szczególności $\E N=G_N'(1-)$,
także gdy średnia jest nieskończona.
Jeżeli $\E N^2<\infty$, to
$\Var(N)=G_N''(1-)+G_N'(1-)-[G_N'(1-)]^2$.
\end{twierdzenie}
\begin{proof}
Promień zbieżności szeregu $\sum_k p_kz^k$ jest co
najmniej jeden, ponieważ $0\leq p_k\leq1$.
Z twierdzenia o różniczkowaniu szeregów potęgowych
wewnątrz koła zbieżności wynika holomorficzność $G_N$
oraz możliwość różniczkowania wyraz po wyrazie:
\[
G_N^{(r)}(z)=\sum_{k=r}^\infty (k)_r p_kz^{k-r}
\quad (|z|<1).
\]
Dla $0<\rho<1$ i $k\geq r$ moduł każdego wyrazu
jest nie większy od $(k)_r\rho^{k-r}$.
Iloraz kolejnych takich ograniczeń wynosi
$\rho(k+1)/(k+1-r)\to\rho<1$, więc szereg ograniczeń
jest zbieżny z kryterium d'Alemberta.
Dla $z=0$ i $r\geq1$ jedynym niezerowym składnikiem
może być wyraz o $k=r$, co daje $G_N^{(r)}(0)=r!p_r$;
dla $r=0$ mamy bezpośrednio $G_N(0)=p_0$.
Dla $0\leq z<1$ definiujemy funkcję na $\N_0$ równą
$(k)_rz^{k-r}$ dla $k\geq r$ i zeru dla $k<r$.
Przedłużając ją zerem na $\R\setminus\N_0$, otrzymujemy
funkcję nieujemną i borelowską. Wzór \eqref{eq:lotus}
oraz zbieżność szeregu pochodnej dają skończoną
wartość oczekiwaną z tezy.
Podstawiając $z=1$ do pierwotnego szeregu, otrzymujemy
$G_N(1)=\sum_kp_k=1$.
Dla $k\geq r$ składnik $(k)_rp_kz^{k-r}$ jest
niemalejący i nieujemny na $[0,1)$, a jego granicą
przy $z\uparrow1$ jest $(k)_rp_k$.
Stosujemy \atlasref{ME-08::thm:mct}{twierdzenie o zbieżności monotonicznej}
względem miary liczącej do dowolnego ciągu $z_j\uparrow1$:
\[
\lim_j G_N^{(r)}(z_j)=
\sum_{k=r}^\infty(k)_rp_k=\E(N)_r.
\]
Monotoniczność $G_N^{(r)}$ na $[0,1)$ gwarantuje istnienie
granicy lewostronnej o tej wartości, także nieskończonej.
Jeżeli $\E N^2<\infty$, to $\E N<\infty$ z nierówności
$N\leq1+N^2$, a $0\leq N(N-1)\leq N^2$.
Wówczas $\E N=G_N'(1-)$, a z tożsamości
$N^2=N(N-1)+N$ otrzymujemy
$\E N^2=G_N''(1-)+G_N'(1-)$.
Odjęcie kwadratu średniej daje podany wzór na wariancję.
\end{proof}

\begin{twierdzenie}[Suma niezależnych liczników]
Jeśli $N_1,N_2$ są niezależnymi zmiennymi losowymi
o wartościach w $\N_0$ na wspólnej przestrzeni
probabilistycznej, to $G_{N_1+N_2}(z)=G_{N_1}(z)G_{N_2}(z)$
dla $|z|\leq1$.
\end{twierdzenie}
\begin{proof}
Oznaczmy $p_j=\Prob(N_1=j)$, $q_k=\Prob(N_2=k)$.
Dla $|z|\leq1$ zmienna $z^{N_1+N_2}$ jest ograniczona
co do modułu przez jeden. Z niezależności
$\Prob(N_1=j,N_2=k)=p_jq_k$.
Stosując \atlasref{ME-13::thm:pushforward-integral-13}{wzór na całkę względem miary obrazowej}
dla wektora $(N_1,N_2)$ oraz wzór na całkę względem
miary dyskretnej, otrzymujemy
\[
\begin{aligned}
G_{N_1+N_2}(z)
&=\sum_{j,k\geq0}z^{j+k}p_jq_k\\
&=\left(\sum_{j\geq0}z^jp_j\right)
  \left(\sum_{k\geq0}z^kq_k\right)
=G_{N_1}(z)G_{N_2}(z).
\end{aligned}
\]
Rozdzielenie sum jest uprawnione, ponieważ
$\sum_{j,k\geq0}|z^{j+k}p_jq_k|\leq
\sum_{j,k\geq0}p_jq_k=1$.
Dla $z=0$ równość $z^{j+k}=z^jz^k$ zachodzi również
przy przyjętej konwencji $0^0=1$.
\end{proof}

\begin{przyklad}[Pięć klasycznych funkcji tworzących]
Korzystamy z definicji rozkładów z rozdziału~\ref{ch:distributions}.
W każdym przypadku $z\in\mathbb C$, $|z|\leq1$.

Dla $B\sim\operatorname{Bern}(p)$, $p\in[0,1]$,
z dwóch mas punktowych otrzymujemy $G_B(z)=1-p+pz$.
Dla $N\sim\Bin(n,p)$, $n\geq0$ całkowitego, z
\eqref{eq:binomial} i dwumianu Newtona wynika
\[
G_N(z)=\sum_{k=0}^n\binom nk (pz)^k(1-p)^{n-k}
=(1-p+pz)^n.
\]
Przy potęgach zerowych stosujemy konwencję $0^0=1$.

Dla $N\sim\Poi(\lambda)$, $\lambda>0$, z szeregu funkcji
wykładniczej dostajemy
\[
G_N(z)=e^{-\lambda}\sum_{k=0}^\infty\frac{(\lambda z)^k}{k!}
=\exp(\lambda(z-1)).
\]
Jeżeli $N_1\sim\Poi(\lambda_1)$ i $N_2\sim\Poi(\lambda_2)$
są niezależne, to ich suma ma funkcję tworzącą
$\exp((\lambda_1+\lambda_2)(z-1))$.
Z jednoznaczności rozkładu wyznaczonego przez funkcję
tworzącą wynika $N_1+N_2\sim\Poi(\lambda_1+\lambda_2)$.
W tym przykładzie $\lambda_1,\lambda_2$ oznaczają parametry
rozkładów, nie miary Lebesgue'a.

Jeśli $T\sim\Geom(p)$, $p\in(0,1]$, oznacza numer pierwszej
udanej próby, to $N=T-1$ liczy poprzedzające ją porażki.
Po oznaczeniu $q=1-p$ mamy $\Prob(N=k)=pq^k$ dla $k\geq0$.
Ponieważ $|qz|\leq q<1$, z szeregu geometrycznego otrzymujemy
\[
G_N(z)=p\sum_{k=0}^\infty(qz)^k=\frac{p}{1-qz}.
\]
W szczególności przy $p=1$ funkcja tworząca jest równa $1$.

Dla liczby porażek $L$ przed $r$-tym sukcesem, gdzie
$r\geq1$ jest liczbą całkowitą, w ciągu niezależnych prób Bernoulliego
z parametrem $p\in(0,1]$, niech $W_j$ oznacza liczbę prób
od poprzedniego sukcesu do następnego, wliczając sukces
(dla $j=1$ liczymy od początku). W rozdziale~\ref{ch:distributions}
wykazano, że $W_1,\ldots,W_r$ są niezależne i mają rozkład
$\Geom(p)$. Zmienne $W_j-1$ również są niezależne:
dla dowolnych borelowskich $A_j$ zdarzenia
$\{W_j-1\in A_j\}=\{W_j\in 1+A_j\}$ są niezależne,
zgodnie z definicją niezależności zmiennych losowych.
Zatem $L=\sum_{j=1}^r(W_j-1)$ prawie na pewno.
Rozkład łączny tych zmiennych jest iloczynem ich rozkładów,
więc wielokrotne sumowanie bezwzględnie zbieżnego szeregu,
jak w dowodzie wzoru dla dwóch niezależnych liczników, daje
\[
G_L(z)=\left(\frac{p}{1-qz}\right)^r.
\]
Dla liczby wszystkich prób $T_r=L+r$ funkcja tworząca
wynosi natomiast $G_{T_r}(z)=z^rG_L(z)$, ponieważ
$z^{L+r}=z^rz^L$, również przy $z=0$ i przyjętej konwencji.
Rozróżnienie tych dwóch konwencji jest istotne.
\end{przyklad}

\section{Funkcja tworząca momenty i kumulanty}

\begin{definicja}
Dla rzeczywistej zmiennej losowej $X$ definiujemy
$M_X(t):=\E e^{tX}\in(0,\infty]$ dla $t\in\R$.
Funkcja podcałkowa jest mierzalna i ściśle dodatnia,
więc jej całka nie może być zerowa: w przeciwnym razie $e^{tX}=0$
prawie na pewno, co jest niemożliwe.
Dziedziną skończoności jest $D_X:=\{t\in\R:M_X(t)<\infty\}$;
zawsze $M_X(0)=1$, zatem $0\in D_X$, lecz $D_X$
nie musi zawierać żadnego otoczenia zera.
Jeśli $M_X$ jest skończona w otoczeniu zera, określamy
funkcję tworzącą kumulanty $K_X(t):=\log M_X(t)$ dla $t\in D_X$.
Kumulantą rzędu $r\geq1$ nazywamy $\kappa_r(X):=K_X^{(r)}(0)$,
gdzie $r$ jest liczbą całkowitą; istnienie tych pochodnych
wynika z następnego twierdzenia.
\end{definicja}

Dziedzina $D_X$ jest przedziałem, być może zredukowanym
do jednego punktu. Istotnie, dla $s,t\in D_X$ i $0<\theta<1$
nierówność Höldera z wykładnikami $1/\theta$ i $1/(1-\theta)$
(\atlasref{ME-14::thm:holder-14}{nierówność Höldera}) daje
\[
M_X(\theta s+(1-\theta)t)
\leq M_X(s)^\theta M_X(t)^{1-\theta}<\infty.
\]
Funkcje $e^{\theta sX}$ oraz $e^{(1-\theta)tX}$ spełniają
hipotezy tej nierówności, ponieważ ich odpowiednie potęgi
mają skończone całki. Przypadki $\theta=0,1$ są bezpośrednie.
Ta sama nierówność po zlogarytmowaniu dowodzi wypukłości
$\log M_X$ na $D_X$.

\begin{twierdzenie}[Istnienie pochodnych i odzyskiwanie momentów]
Niech $X$ będzie rzeczywistą zmienną losową oraz
$\delta>0$ takie, że $M_X(-\delta)<\infty$ i
$M_X(\delta)<\infty$. Wtedy $M_X$ jest klasy $C^\infty$ na $(-\delta,\delta)$ oraz
dla każdego $r\in\N_0$ i $|t|<\delta$ zachodzi
\[
M_X^{(r)}(t)=\E[X^r e^{tX}].
\]
W szczególności $\E|X|^r<\infty$ i
$M_X^{(r)}(0)=\E X^r$ dla każdego $r\in\N_0$.
Ponadto dla $|t|<\delta$ zachodzi rozwinięcie
\[
M_X(t)=\sum_{r=0}^\infty\frac{\E X^r}{r!}t^r.
\]
\end{twierdzenie}
\begin{proof}
Najpierw zwróćmy uwagę, że
$e^{\delta|X|}\leq e^{\delta X}+e^{-\delta X}$.
Zatem $\E e^{\delta|X|}<\infty$.
Ustalmy $0<\eta<\delta$ i całkowite $r\geq0$.
Funkcja $u\mapsto u^re^{-(\delta-\eta)u}$ jest ograniczona na
$[0,\infty)$, ponieważ dla $r>0$ jej granica w nieskończoności
jest zerowa, a dla $r=0$ nie przekracza jedności.
Istnieje zatem stała $C_{r,\eta,\delta}<\infty$ taka, że
dla $|t|\leq\eta$:
\[
|X|^re^{tX}\leq C_{r,\eta,\delta}e^{\delta|X|}.
\]
W szczególności wszystkie momenty bezwzględne rzędów
całkowitych są skończone (weźmy $t=0$), a skończoność
momentów pozostałych dodatnich rzędów wynika z oszacowania
momentów niższych rzędów z poprzedniej sekcji. Oznaczmy $H_r(t)=\E[X^re^{tX}]$,
przyjmując $X^0=1$. Funkcja $H_r$ jest ciągła na
$(-\delta,\delta)$ przez
\atlasref{ME-08::thm:dct}{twierdzenie o zbieżności zmajoryzowanej}
i powyższą wspólną majorantę na każdym przedziale
$[-\eta,\eta]$.

Dla ustalonego $t$ wybieramy $\eta$ tak, aby
$|t|<\eta<\delta$. Jeżeli $0<|h|<\eta-|t|$, to
z twierdzenia o wartości średniej dla funkcji wykładniczej
dla każdego $x\in\R$ mamy
\[
\left|x^r\frac{e^{(t+h)x}-e^{tx}}{h}\right|
\leq |x|^{r+1}e^{\eta|x|}
\leq C_{r+1,\eta,\delta}e^{\delta|x|}.
\]
Iloraz różnicowy dąży punktowo do $x^{r+1}e^{tx}$.
Zbieżność zmajoryzowana pozwala więc przenieść granicę
pod całkę i daje $H_r'(t)=H_{r+1}(t)$.
Ponieważ $M_X=H_0$, indukcja dowodzi
$M_X^{(r)}=H_r$, a ciągłość wszystkich $H_r$ daje klasę
$C^\infty$. Dla $t=0$ otrzymujemy żądane wzory na momenty.
Aby uzyskać rozwinięcie w szereg, dla ustalonego $|t|<\delta$
stosujemy szereg wykładniczy. Moduł każdej sumy częściowej
$\sum_{r=0}^n(tX)^r/r!$ nie przekracza $e^{|t||X|}$,
a ta funkcja jest całkowalna, ponieważ
$e^{|t||X|}\leq e^{\delta|X|}$. Zbieżność zmajoryzowana
daje wskazany szereg dla $M_X(t)$. Jego bezwzględna zbieżność
wynika z oszacowania
\[
\sum_{r=0}^\infty\frac{|\E X^r|}{r!}|t|^r
\leq \E\left[\sum_{r=0}^\infty\frac{|tX|^r}{r!}\right]
=\E e^{|t||X|}<\infty,
\]
gdzie wymiana sumy z całką wynika ze zbieżności monotonicznej.
\end{proof}

\begin{twierdzenie}[Dodawanie kumulantów]
Niech $X,Y$ będą niezależnymi rzeczywistymi zmiennymi losowymi
na wspólnej przestrzeni probabilistycznej, a ich funkcje $M_X,M_Y$ będą skończone w pewnym wspólnym
otoczeniu zera. Wtedy w tym otoczeniu
$M_{X+Y}(t)=M_X(t)M_Y(t)$ i
$K_{X+Y}(t)=K_X(t)+K_Y(t)$.
W szczególności $\kappa_r(X+Y)=\kappa_r(X)+\kappa_r(Y)$
dla każdego całkowitego $r\geq1$.
Ponadto $K_X'(0)=\E X$, $K_X''(0)=\Var(X)$
oraz $K_X'''(0)=\E(X-\E X)^3$.
\end{twierdzenie}
\begin{proof}
Z niezależności rozkład łączny jest miarą produktową
$\mu_X\otimes\mu_Y$ (rozdział~\ref{ch:random-variables}).
Korzystając z całkowania względem miary obrazowej i
\atlasref{ME-12::thm:tonelli-12}{twierdzenia Tonellego}
dla nieujemnej borelowskiej funkcji $(x,y)\mapsto e^{tx}e^{ty}$,
otrzymujemy
\[
M_{X+Y}(t)=\int_{\R^2}e^{tx}e^{ty}\,
\dd(\mu_X\otimes\mu_Y)(x,y)
=M_X(t)M_Y(t).
\]
Oba czynniki są dodatnie i skończone we wspólnym otoczeniu
zera, zatem $M_{X+Y}$ jest tam skończona, a logarytmowanie
daje wzór na sumę funkcji tworzących kumulanty.
Z założenia o wspólnym otoczeniu zera można wybrać
$\delta_0>0$ tak, aby wszystkie trzy funkcje tworzące
momenty były skończone w punktach $\pm\delta_0$.
Poprzednie twierdzenie i gładkość logarytmu na $(0,\infty)$
zapewniają istnienie wszystkich pochodnych kumulant.
Różniczkując tę równość w zerze, uzyskujemy
$\kappa_r(X+Y)=\kappa_r(X)+\kappa_r(Y)$ dla każdego $r\geq1$. Ponieważ $M_X(0)=1$, reguły
różniczkowania logarytmu dają
$K_X'=M_X'/M_X$ i
$K_X''=M_X''/M_X-(M_X'/M_X)^2$.
W szczególności $K_X'(0)=\E X$, a
$K_X''(0)=\E X^2-(\E X)^2=\Var(X)$.
Różniczkując jeszcze raz, dostajemy
\[
K_X'''=\frac{M_X'''}{M_X}
-\frac{3M_X'M_X''}{M_X^2}
+\frac{2(M_X')^3}{M_X^3}.
\]
Podstawiając $t=0$ i wzory na momenty, otrzymujemy
$K_X'''(0)=\E X^3-3\E X\,\E X^2+2(\E X)^3$.
Rozwinięcie $(X-\E X)^3$ i liniowość całki pokazują,
że jest to trzeci moment centralny.
\end{proof}

\begin{przyklad}[Kumulanty trzech rodzin]
Korzystamy ze wzorów na funkcje tworzące momenty
z rozdziału~\ref{ch:distributions}.
Dla $X\sim\Normal(\mu,\sigma^2)$, $\mu\in\R$, $\sigma>0$,
logarytm funkcji tworzącej daje
$K_X(t)=\mu t+\sigma^2t^2/2$, zatem
$\kappa_1(X)=\mu$, $\kappa_2(X)=\sigma^2$
i $\kappa_r(X)=0$ dla $r\geq3$.
Dla $N\sim\Poi(\lambda)$, $\lambda>0$, zachodzi
$K_N(t)=\lambda(e^t-1)$, a więc
$\kappa_r(N)=\lambda$ dla każdego $r\geq1$.
Dla $B\sim\operatorname{Bern}(p)$, $p\in[0,1]$,
mamy $M_B(t)=1-p+pe^t$ oraz $K_B(t)=\log(1-p+pe^t)$.
Ponieważ $B^r=B$ dla całkowitych $r\geq1$, poprzednie
wzory dają $\kappa_1(B)=p$, $\kappa_2(B)=p(1-p)$
oraz $\kappa_3(B)=p-3p^2+2p^3=p(1-p)(1-2p)$.
\end{przyklad}

\begin{stwierdzenie}[Jednoznaczność funkcji tworzącej momenty]
Niech $X,Y$ będą rzeczywistymi zmiennymi losowymi.
Jeśli ich funkcje $M_X,M_Y$ są skończone na
$(-\delta,\delta)$ dla pewnego $\delta>0$
i są tam równe, to $\Law(X)=\Law(Y)$.
\end{stwierdzenie}
\begin{proof}
Rozważmy transformaty zespolone
$L_X(z)=\E e^{zX}$ i $L_Y(z)=\E e^{zY}$
na pasie $\{z\in\mathbb C:|\operatorname{Re}z|<\delta\}$.
Wykażemy holomorficzność $L_X$; dla $L_Y$ argument
jest identyczny. Dla ustalonego $z_0$ z pasa wybierzmy
$|\operatorname{Re}z_0|<\eta<\delta'<\delta$.
Mamy $\E e^{\delta'|X|}<\infty$ z oszacowania
$e^{\delta'|X|}\leq e^{\delta'X}+e^{-\delta'X}$,
więc całka definiująca $L_X(z_0)$ jest bezwzględnie zbieżna.
Dla zespolonego $w\ne0$, $|w|<\eta-|\operatorname{Re}z_0|$,
korzystamy z tożsamości
\[
\frac{e^{(z_0+w)x}-e^{z_0x}}{w}
=x\int_0^1 e^{(z_0+sw)x}\,\dd\lambda_1(s).
\]
Wynika ona z podstawowego twierdzenia rachunku całkowego
dla funkcji zmiennej rzeczywistej $s$, zastosowanego
do części rzeczywistej i urojonej.
Moduł prawej strony nie przekracza
$|x|e^{\eta|x|}\leq C e^{\delta'|x|}$.
Iloraz dąży do $xe^{z_0x}$ przy $w\to0$.
Zbieżność zmajoryzowana, stosowana do obu części,
daje istnienie pochodnej zespolonej
$L_X'(z_0)=\E[Xe^{z_0X}]$.
Ponieważ $z_0$ jest dowolny, $L_X$ jest holomorficzna
na całym pasie.
Pas jest otwartym i spójnym zbiorem, a przedział
$(-\delta,\delta)$ ma w nim punkt skupienia.
Z zasady identyczności funkcji holomorficznych równość
$L_X(t)=L_Y(t)$ na tym przedziale implikuje
równość na całym pasie, w szczególności dla $z=it$.
Zatem $\varphi_X(t)=\varphi_Y(t)$ dla każdego $t\in\R$.
Twierdzenie o jednoznaczności funkcji charakterystycznej,
udowodnione w następnej sekcji, daje $\Law(X)=\Law(Y)$.
\end{proof}

W szczególności jeżeli obie funkcje tworzące momenty
są skończone w pewnym wspólnym otoczeniu zera i
$\E X^r=\E Y^r$ dla wszystkich całkowitych $r\geq1$,
to rozwinięcia w szeregi potęgowe są identyczne
na dostatecznie małym otoczeniu zera (wyrazy stałe
są równe jeden). Z udowodnionego stwierdzenia wynika
wtedy $\Law(X)=\Law(Y)$. Samo istnienie wszystkich
momentów nie zapewnia jednak takiej jednoznaczności,
co pokazuje następny przykład.

\begin{przyklad}[Wszystkie momenty nie muszą wyznaczać rozkładu]
Niech $Z\sim\Normal(0,1)$ oraz $X=e^Z$.
Ze wzoru na funkcję tworzącą standardowego rozkładu normalnego
(rozdział~\ref{ch:distributions})
otrzymujemy $\E X^n=\E e^{nZ}=e^{n^2/2}$ dla każdego
całkowitego $n\geq0$. Jednak dla $t>0$
\[
M_X(t)=\frac{1}{\sqrt{2\pi}}\int_\R
\exp(te^z-z^2/2)\,\dd\lambda_1(z)=\infty.
\]
Istotnie, wykładnik $te^z-z^2/2$ dąży do $+\infty$
przy $z\to+\infty$, więc funkcja podcałkowa bez czynnika
$1/\sqrt{2\pi}$ jest od pewnego miejsca nie mniejsza od $1$.
Dla $t\leq0$ mamy natomiast $0<e^{tX}\leq1$,
zatem $D_X=(-\infty,0]$.

Z twierdzenia o transformacji monotonicznej dla
$\exp:\R\to(0,\infty)$ wynika, że $X$ ma gęstość
\[
f_0(x)=
\begin{cases}
\displaystyle\frac{\exp(-(\log x)^2/2)}{x\sqrt{2\pi}},&x>0,\\
0,&x\leq0.
\end{cases}
\]
Dla $\varepsilon\in[-1,1]$ połóżmy
$f_\varepsilon(x)=f_0(x)[1+\varepsilon\sin(2\pi\log x)]$
dla $x>0$, a $f_\varepsilon(x)=0$ dla $x\leq0$.
Funkcje te są borelowskie, nieujemne i ograniczone
z góry przez $2f_0$. W szczególności ich całki i wszystkie
momenty nieujemnych rzędów całkowitych są skończone.
Aby sprawdzić normalizację i równość momentów, podstawiamy
$x=e^z$ w całce. Dla każdego całkowitego $n\geq0$
\[
\begin{aligned}
\int_0^\infty x^n f_0(x)\sin(2\pi\log x)\,
\dd\lambda_1(x)
&=\operatorname{Im}\E e^{(n+2\pi i)Z}\\
&=\operatorname{Im}\exp((n+2\pi i)^2/2)=0.
\end{aligned}
\]
Ostatni wzór wynika z holomorficzności $L_Z$ na
każdym pasie $|\operatorname{Re}z|<\delta$, wykazanej
w poprzednim dowodzie, oraz zasady identyczności:
$L_Z(z)=e^{z^2/2}$ na $\mathbb C$, ponieważ obie funkcje
są holomorficzne i zgadzają się na $\R$.
Część urojona wynosi $e^{n^2/2-2\pi^2}\sin(2\pi n)=0$.
Dla $n=0$ otrzymujemy $\int_\R f_\varepsilon\,\dd\lambda_1=1$,
a dla wszystkich $n\geq1$ identyczne momenty $e^{n^2/2}$.
Dla $\varepsilon\ne\varepsilon'$ funkcje
$f_\varepsilon$ i $f_{\varepsilon'}$ różnią się na
$(0,\infty)$ poza przeliczalnym zbiorem
$\{e^{k/2}:k\in\mathbb Z\}$, gdyż $f_0$ jest tam dodatnia.
Z jednoznaczności gęstości prawie wszędzie wynika,
że wyznaczają różne rozkłady. Ten przykład pokazuje, dlaczego sama lista momentów
bez warunków zapewniających jednoznaczność może być niewystarczająca.
\end{przyklad}

\section{Funkcja charakterystyczna i transformata Fouriera}

Kluczowa różnica między $e^{tX}$ a $e^{itX}$ polega na tym,
że $|e^{itX}|=1$ nawet dla bardzo ciężkich ogonów rozkładu.
Całkowanie zespolonej wykładniczej nie wymaga więc istnienia
żadnego momentu. Zmienna faza $e^{itx}$ rejestruje zarazem
informację o całym rozkładzie, nie tylko o jego średniej.

\begin{definicja}
Funkcją charakterystyczną rzeczywistej zmiennej losowej $X$
nazywamy $\varphi_X:\R\to\mathbb C$ zadaną wzorem
$\varphi_X(t):=\E e^{itX}=\int_\R e^{itx}\,\dd\mu_X(x)$.
Równość wynika z \eqref{eq:lotus}, ponieważ funkcja
$x\mapsto e^{itx}$ jest ciągła i ma moduł jeden.
Dla borelowskiej miary probabilistycznej $\mu$ piszemy
analogicznie $\varphi_\mu(t)=\int_\R e^{itx}\,\dd\mu(x)$.
Jest to transformata Fouriera miary z konwencją znaku
dodatniego w wykładniku.
\end{definicja}

\begin{twierdzenie}[Własności funkcji charakterystycznej]
Dla każdej rzeczywistej zmiennej losowej $X$ i
$t,a,b\in\R$ zachodzą:
\begin{enumerate}
\item $\varphi_X(0)=1$, $|\varphi_X(t)|\leq1$
 i $\varphi_X(-t)=\overline{\varphi_X(t)}$;
\item $\varphi_X$ jest jednostajnie ciągła na $\R$;
\item $\varphi_{aX+b}(t)=e^{itb}\varphi_X(at)$;
\item dla niezależnych rzeczywistych zmiennych losowych
 $X,Y$ na wspólnej przestrzeni probabilistycznej mamy
 $\varphi_{X+Y}(t)=\varphi_X(t)\varphi_Y(t)$;
\item dla dowolnych $n\in\N$, $t_1,\ldots,t_n\in\R$,
 $c_1,\ldots,c_n\in\mathbb C$ zachodzi
 $\sum_{j,k=1}^n c_j\overline{c_k}\varphi_X(t_j-t_k)\geq0$.
\end{enumerate}
Jeżeli $\E|X|^r<\infty$ dla pewnego całkowitego $r\geq1$,
to $\varphi_X$ jest klasy $C^r$ na $\R$ i
$\varphi_X^{(k)}(t)=\E[(iX)^k e^{itX}]$
dla $0\leq k\leq r$.
\end{twierdzenie}
\begin{proof}
Ponieważ $|e^{itX}|=1$, całka zawsze istnieje,
a nierówność trójkąta daje $|\varphi_X(t)|\leq1$.
Pozostałe tożsamości z punktu pierwszego wynikają
z podstawienia $t=0$ i sprzężenia całki.
Dla ciągu $h_n\to0$ wyrażenie
$|\varphi_X(t+h_n)-\varphi_X(t)|$
jest nie większe niż $\E|e^{ih_nX}-1|$, niezależnie od $t$.
Funkcja pod wartością oczekiwaną dąży punktowo do zera
i nie przekracza $2$. Z
\atlasref{ME-08::thm:dct}{twierdzenia o zbieżności zmajoryzowanej}
wynika $\E|e^{ihX}-1|\to0$ przy $h\to0$.
Ponieważ oszacowanie jest niezależne od $t$, otrzymujemy
$\sup_{t\in\R}|\varphi_X(t+h)-\varphi_X(t)|\to0$,
czyli jednostajną ciągłość. Punkt trzeci wynika z równości
$e^{it(aX+b)}=e^{itb}e^{i(at)X}$.
Dla czwartego korzystamy z niezależności, dzięki której
rozkład łączny jest miarą $\mu_X\otimes\mu_Y$.
Funkcja $(x,y)\mapsto e^{itx}e^{ity}$ ma moduł jeden,
jest więc całkowalna względem tej miary probabilistycznej.
Z twierdzenia o mierze obrazowej i
\atlasref{ME-12::thm:fubini-12}{twierdzenia Fubiniego}
dla obu części zespolonych wynika
$\E[e^{itX}e^{itY}]=\E e^{itX}\E e^{itY}$.
Dla ostatniego punktu korzystamy z liniowości całki
zespolonej dla skończonej sumy:
\[
\begin{aligned}
\sum_{j,k=1}^n c_j\overline{c_k}\varphi_X(t_j-t_k)
&=\E\left[\sum_{j,k=1}^n
c_je^{it_jX}\overline{c_ke^{it_kX}}\right]\\
&=\E\left|\sum_{j=1}^n c_je^{it_jX}\right|^2\geq0.
\end{aligned}
\]
Ostatnia funkcja jest nieujemna i ograniczona przez
$(\sum_{j=1}^n|c_j|)^2$, a więc całkowalna.
Wreszcie oznaczmy $H_k(t)=\E[(iX)^ke^{itX}]$,
$0\leq k\leq r$, przy konwencji $X^0=1$.
Z nierówności $|X|^k\leq1+|X|^r$ wynika całkowalność
tych funkcji podcałkowych. Zbieżność zmajoryzowana
daje ciągłość wszystkich $H_k$.
Dla $0\leq k<r$ i rzeczywistego $h\ne0$ zachodzi
\[
\left|(ix)^k e^{itx}\frac{e^{ihx}-1}{h}\right|
\leq |x|^{k+1},
\]
gdyż $e^{ihx}-1=ihx\int_0^1e^{ishx}\,\dd\lambda_1(s)$.
Iloraz różnicowy dąży do $(ix)^{k+1}e^{itx}$,
a majoranta $|X|^{k+1}$ jest całkowalna.
Zbieżność zmajoryzowana zastosowana do obu części
zespolonych daje $H_k'=H_{k+1}$.
Indukcja od $H_0=\varphi_X$ oraz ciągłość $H_r$
dowodzą tezy.
\end{proof}

\begin{lemat}[Całka Fouriera dla rozkładu Gaussa]
Jeśli $Z\sim\Normal(0,1)$, to
$\varphi_Z(t)=e^{-t^2/2}$ dla $t\in\R$.
Jeśli $Y\sim\Normal(m,\sigma^2)$, gdzie $m\in\R$ i
$\sigma>0$, to
$\varphi_Y(t)=\exp(imt-\sigma^2t^2/2)$.
\end{lemat}
\begin{proof}
Niech $g(x)=(2\pi)^{-1/2}e^{-x^2/2}$.
Z definicji gęstości normalnej (rozdział~\ref{ch:distributions})
oraz \eqref{eq:lotus} mamy
$\E|Z|=2(2\pi)^{-1/2}\int_0^\infty xe^{-x^2/2}\,\dd\lambda_1(x)
=\sqrt{2/\pi}<\infty$, gdyż podstawienie $u=x^2/2$
daje $\int_0^\infty e^{-u}\,\dd\lambda_1(u)=1$.
Poprzednie twierdzenie daje więc
$\varphi_Z'(t)=i\int_\R xe^{itx}g(x)\,\dd\lambda_1(x)$.
Ponieważ $g'(x)=-xg(x)$, całkowanie przez części na
$[-R,R]$ (osobno dla części rzeczywistej i urojonej) daje
\[
\int_{-R}^R xe^{itx}g(x)\,\dd\lambda_1(x)
=-[e^{itx}g(x)]_{-R}^R+
it\int_{-R}^R e^{itx}g(x)\,\dd\lambda_1(x).
\]
Obie całki przechodzą do całek po $\R$ przy $R\to\infty$,
ponieważ odpowiednio $|x|g(x)$ i $g(x)$ są całkowalne.
Wyraz brzegowy ma moduł nie większy od $2g(R)\to0$.
Stąd $\varphi_Z'(t)=-t\varphi_Z(t)$, czyli
$(e^{t^2/2}\varphi_Z(t))'=0$.
Warunek $\varphi_Z(0)=1$ daje $\varphi_Z(t)=e^{-t^2/2}$.
Z transformacji gęstości przez $x\mapsto m+\sigma x$
wynika $\Law(Y)=\Law(m+\sigma Z)$.
Własność afiniczna funkcji charakterystycznej daje drugą równość.
\end{proof}

\begin{twierdzenie}[Jednoznaczność funkcji charakterystycznej]
Jeżeli $X,Y$ są rzeczywistymi zmiennymi losowymi i
$\varphi_X(t)=\varphi_Y(t)$ dla każdego $t\in\R$,
to $\Law(X)=\Law(Y)$.
\end{twierdzenie}
\begin{proof}
Oznaczmy rozkłady przez $\mu$ i $\nu$, a przez
$g_\sigma(x)=(\sigma\sqrt{2\pi})^{-1}
\exp(-x^2/(2\sigma^2))$ gęstość $\Normal(0,\sigma^2)$
dla $\sigma>0$. Oznaczmy przez $\gamma_\sigma$ miarę
$g_\sigma\lambda_1$ i przez $\mu*\gamma_\sigma$ obraz
miary produktowej $\mu\otimes\gamma_\sigma$ przez dodawanie.
Dla zbioru borelowskiego $B$ korzystamy z definicji
miary obrazowej, podstawienia $x=y+u$ (przesunięcie
zachowuje miarę Lebesgue'a) i twierdzenia Tonellego
dla nieujemnej borelowskiej funkcji
$(y,u)\mapsto\chi_B(y+u)g_\sigma(u)$. Dostajemy
\[
(\mu*\gamma_\sigma)(B)
=\int_\R\int_\R\chi_B(y+u)g_\sigma(u)\,
\dd\lambda_1(u)\,\dd\mu(y)
=\int_B\int_\R g_\sigma(x-y)\,\dd\mu(y)\,\dd\lambda_1(x).
\]
Zatem funkcja $x\mapsto\int_\R g_\sigma(x-y)\,\dd\mu(y)$
jest gęstością splotu; jej mierzalność wynika z Tonellego.
Z poprzedniego lematu, podstawiając $v=\sigma t$, otrzymujemy
dla każdego $u\in\R$
\[
\frac1{2\pi}\int_\R e^{-itu}e^{-\sigma^2t^2/2}\,\dd\lambda_1(t)
=\frac1{\sigma\sqrt{2\pi}}e^{-u^2/(2\sigma^2)}
=g_\sigma(u).
\]
Całka z modułu funkcji podcałkowej po $\mu(y)$ i
$\lambda_1(t)$ wynosi $\int_\R e^{-\sigma^2t^2/2}\,\dd\lambda_1(t)
=\sqrt{2\pi}/\sigma<\infty$.
Możemy więc zastosować Fubiniego i dla każdego $x\in\R$
otrzymujemy
\[
\int_\R g_\sigma(x-y)\,\dd\mu(y)
=\frac1{2\pi}\int_\R
 e^{-itx}\varphi_\mu(t)e^{-\sigma^2t^2/2}
\,\dd\lambda_1(t).
\]
To samo zachodzi dla $\nu$.
Równość funkcji charakterystycznych oznacza zatem,
że sploty $\mu*\gamma_\sigma$ i $\nu*\gamma_\sigma$ mają
identyczne gęstości, a więc są równymi miarami.

Dla dowolnej funkcji $h:\R\to\R$ ograniczonej i ciągłej
oraz zmiennej $Z\sim\Normal(0,1)$, określonej na dowolnej
przestrzeni probabilistycznej, zachodzi
\[
\int_\R h(x)\,\dd(\mu*\gamma_\sigma)(x)
=\int_\R\E[h(y+\sigma Z)]\,\dd\mu(y).
\]
Jest to wzór na całkę względem miary obrazowej i Fubini
dla funkcji ograniczonej na przestrzeni o mierze jeden;
$\sigma Z$ ma rozkład $\gamma_\sigma$. Wewnętrzna
wartość oczekiwana jest borelowską funkcją $y$,
co wynika z mierzalności całki względem parametru
w twierdzeniu Fubiniego (lub z ciągłości i zbieżności
zmajoryzowanej).
Gdy $\sigma\downarrow0$, dla każdego ustalonego $y$
zmienna $h(y+\sigma Z)$ dąży punktowo do $h(y)$,
a jej moduł jest ograniczony przez $\|h\|_\infty$.
Dwukrotne zastosowanie zbieżności zmajoryzowanej daje
$\int h\,\dd(\mu*\gamma_\sigma)\to\int h\,\dd\mu$;
podobnie dla $\nu$. Ponieważ sploty były równe dla
każdego $\sigma$, otrzymujemy
$\int h\,\dd\mu=\int h\,\dd\nu$ dla każdej
ograniczonej funkcji ciągłej $h$. Dla ustalonego $a\in\R$ niech $h_n$ będzie równa $1$
na $(-\infty,a]$, liniowo maleje od $1$ do $0$ na
$[a,a+1/n]$ i wynosi $0$ na $[a+1/n,\infty)$.
Funkcje te są ograniczone, ciągłe i dążą punktowo
do $\chi_{(-\infty,a]}$. Zbieżność zmajoryzowana
względem obu miar daje $\mu((-\infty,a])=\nu((-\infty,a])$.
Z jednoznaczności miary wyznaczonej przez dystrybuantę
(rozdział~\ref{ch:random-variables}) otrzymujemy $\mu=\nu$.
\end{proof}

\begin{stwierdzenie}[Odwracanie transformaty przy całkowalności]
Niech $X$ będzie rzeczywistą zmienną losową taką,
że $\int_\R|\varphi_X(t)|\,\dd\lambda_1(t)<\infty$.
Wtedy rozkład $X$ ma ograniczoną, ciągłą gęstość
\[
f_X(x)=\frac1{2\pi}\int_\R e^{-itx}\varphi_X(t)
\,\dd\lambda_1(t),\qquad x\in\R.
\]
\end{stwierdzenie}
\begin{proof}
Niech $\gamma_\sigma=g_\sigma\lambda_1$ będzie miarą
gaussowską z poprzedniego dowodu. Splot $\mu_X*\gamma_\sigma$
ma gęstość
$f_\sigma(x)=(2\pi)^{-1}\int_\R e^{-itx}\varphi_X(t)
 e^{-\sigma^2t^2/2}\,\dd\lambda_1(t)$.
Oznaczmy przez $f$ funkcję z tezy. Dla każdego $x$
\[
|f_\sigma(x)-f(x)|\leq\frac1{2\pi}\int_\R
|\varphi_X(t)|\,|e^{-\sigma^2t^2/2}-1|\,\dd\lambda_1(t).
\]
Prawa strona jest niezależna od $x$ i dąży do zera
przez zbieżność zmajoryzowaną, z majorantą $|\varphi_X|$.
Zatem $f_\sigma\to f$ jednostajnie na $\R$, a w szczególności
$f$ jest rzeczywista i nieujemna, gdyż takie są $f_\sigma$.
Z nierówności trójkąta mamy $|f(x)|\leq
(2\pi)^{-1}\int_\R|\varphi_X(t)|\,\dd\lambda_1(t)$.
Ciągłość $f$ wynika ze zbieżności zmajoryzowanej przy
$x_n\to x$ w całce definiującej $f$.
Dla nieujemnej funkcji ciągłej $h$ o zwartym nośniku
mamy
\[
\left|\int_\R h(x)(f_\sigma(x)-f(x))\,\dd\lambda_1(x)\right|
\leq \|f_\sigma-f\|_\infty\int_\R h(x)\,\dd\lambda_1(x)
\longrightarrow0,
\]
gdyż $h$ jest ograniczona i ma zwarty nośnik,
a jednocześnie $\int_\R h(x)f_\sigma(x)\,\dd\lambda_1(x)
\to\E h(X)$
z argumentu z dowodu jednoznaczności.
Zatem $\int_\R h(x)f(x)\,\dd\lambda_1(x)=\E h(X)$
dla wszystkich takich $h$.
Aby uzyskać normalizację, weźmy
$h_n(x)=\min(1,\max(0,n+1-|x|))$, $n\geq1$.
Są to nieujemne ciągłe funkcje o zwartym nośniku
i $h_n\uparrow1$. Zbieżność monotoniczna po obu stronach
daje $\int_\R f(x)\,\dd\lambda_1(x)=1$.
Miara $\rho=f\lambda_1$ jest więc probabilistyczna.
Dla dowolnej ograniczonej ciągłej funkcji $h\geq0$ stosujemy
uzyskaną równość do $hh_n$ i przechodzimy do granicy
przez zbieżność monotoniczną. Otrzymujemy
$\int_\R h\,\dd\rho=\int_\R h\,\dd\mu_X$.
Stosując ciągłą aproksymację półprostych z poprzedniego
dowodu, wnioskujemy $\rho=\mu_X$.
\end{proof}

\begin{przyklad}[Ciężkie ogony a funkcja charakterystyczna]
Dla standardowego rozkładu Cauchy'ego zachodzi
$\varphi_X(t)=e^{-|t|}$. Udowodnimy tę tożsamość
bez odwoływania się do rachunku residuów.
Niech $Z_1,Z_2$ będą niezależnymi zmiennymi standardowo
normalnymi. Z przedstawienia rozkładu Studenta dla
jednego stopnia swobody (rozdział~\ref{ch:distributions})
wiemy, że $C=Z_1/|Z_2|$ ma rozkład Cauchy'ego;
na zdarzeniu $\{Z_2=0\}$, którego prawdopodobieństwo
wynosi zero, definiujemy $C=0$.
Tak określona zmienna $C$ jest mierzalna: iloraz jest
ciągły na otwartym zbiorze par $(x,z)$ z $z\ne0$,
a na domkniętym zbiorze $\{z=0\}$ przyjmujemy zero.
Rozkład łączny $Z_1,Z_2$ jest miarą produktową.
Funkcja $e^{itC}$ ma moduł jeden, więc po przeniesieniu
całkowania na rozkład łączny Fubini pozwala
całkować najpierw po $Z_1$ dla ustalonego $Z_2=z\ne0$.
Wewnętrzna całka to $\varphi_{Z_1}(t/|z|)$, czyli
$e^{-t^2/(2z^2)}$. Korzystając z symetrii gęstości $Z_2$,
otrzymujemy poniższą równość, przyjmując funkcję
$x\mapsto\exp(-t^2/(2x^2))$ równą $1$ w punkcie $x=0$
(co odpowiada przyjętemu $C=0$ na $\{Z_2=0\}$):
\[
\varphi_C(t)=\E\exp\!\left(-\frac{t^2}{2Z_2^2}\right)
=\sqrt{\frac2\pi}\int_0^\infty
\exp\!\left(-\frac{z^2}{2}-\frac{t^2}{2z^2}\right)
\,\dd\lambda_1(z).
\]
Aby obliczyć tę całkę, dla $a>0$ i $b\geq0$ połóżmy
$I(a,b)=\int_0^\infty e^{-az^2-b/z^2}\,\dd\lambda_1(z)$.
Całka jest dodatnia i skończona, ponieważ jej funkcja
podcałkowa jest dodatnia i nie większa od $e^{-az^2}$.
Dla $b>0$ ustalmy $0<b_0<b<b_1$. Z twierdzenia o wartości
średniej ilorazy różnicowe względem $b$, dla przyrostów
pozostających w $[b_0,b_1]$, są ograniczone przez
$z^{-2}e^{-az^2-b_0/z^2}$.
Na $(0,1]$ całkowalność tej majoranty wynika z podstawienia
$v=1/z$, gdyż $\int_1^\infty e^{-b_0v^2}\,\dd\lambda_1(v)<\infty$;
na $[1,\infty)$ majoranta nie przekracza $e^{-az^2}$.
Zbieżność zmajoryzowana uzasadnia różniczkowanie.
Podstawienie $u=\sqrt{b/a}/z$ w całce pochodnej daje
$z^{-2}|\dd z/\dd u|=\sqrt{a/b}$ (gdyż
$z=\sqrt{b/a}/u$ i $|\dd z/\dd u|=\sqrt{b/a}/u^2$) oraz
$az^2+b/z^2=b/u^2+au^2$. Stąd
\[
\frac{\partial I}{\partial b}(a,b)
=-\int_0^\infty z^{-2}e^{-az^2-b/z^2}\,\dd\lambda_1(z)
=-\sqrt{\frac ab}\,I(a,b).
\]
Zatem dla $b>0$ mamy $\partial_b\log I(a,b)=-\sqrt{a/b}$.
Pochodna funkcji $b\mapsto\log I(a,b)+2\sqrt{ab}$
jest zerowa, więc funkcja ta jest stała na $(0,\infty)$.
Po potęgowaniu otrzymujemy $I(a,b)=C(a)e^{-2\sqrt{ab}}$.
Zbieżność zmajoryzowana przy $b\downarrow0$ z majorantą
$e^{-az^2}$ daje $C(a)=I(a,0)$.
Podstawiając $u=\sqrt{2a}\,z$ i korzystając z całki Gaussa
(rozdział~\ref{ch:distributions}), otrzymujemy
$I(a,0)=\sqrt\pi/(2\sqrt a)$.
Podstawiając $a=1/2$ i $b=t^2/2$ do powyższego
wzoru na $\varphi_C$, otrzymujemy
$\varphi_C(t)=e^{-|t|}$ również dla $t=0$.
Funkcja charakterystyczna istnieje więc pomimo braku
wartości oczekiwanej rozkładu Cauchy'ego.
\end{przyklad}

